\section{Weighted enumeration and grouped coefficient formulas}
\label{sec:incremental}

The coefficient theorem admits several finite constructions at a prescribed cutoff. This section organizes them by exact weight, by increasing Newton precision, and by marker degree. All constructions use finitely many positive gaps $\delta_i$ and preserve the meaning of a complete power--log block.

\subsection{Enumeration by complete weight layers}

Let $I$ be the set of retained multi-indices, initially empty, and let its
candidate frontier initially consist of the zero multi-index. At each step,
select the smallest candidate weight $\sigma$, promote \emph{every} candidate
of that weight into $I$, and add all their one-step successors
$\bk+e_i$ to the candidate set. The frontier is a set of integer vectors, so a candidate reached through several predecessors occurs only once. Distinct candidates of equal weight remain distinct until their coefficients are summed.

\begin{proposition}[Complete layers and coefficient reuse]
After the layer at weight $\sigma$ is processed, $I$ is exactly
$\{\bk\in\N^m:\bk\cdot\bdelta\le\sigma\}$, and the candidate set is its
one-step outer boundary. This construction determines every retained multi-index coefficient once. Summing the coefficients of an entire layer before
counting a block preserves all resonances and exact cancellations.
\end{proposition}
\begin{proof}
Every nonzero multi-index has a predecessor of strictly smaller weight because
all gaps are positive. Inductively, every index of the next candidate weight
has already been generated by a processed predecessor. A successor of an index
in the current layer has strictly larger weight, so processing one member of
the layer cannot expose a previously unseen member of the same layer.
The set union identifies repeated occurrences of a candidate through different predecessors. The
coefficient of each weight is the complete finite sum in
Theorem~\ref{thm:coefficients}, whether or not that sum is zero.
\end{proof}

Extending a cutoff requires only subsequent layers: coefficients below the
previous cutoff do not depend on the new cutoff. A prescribed exponent
cutoff is reached in finitely many layers by local finiteness. A prescribed
number of nonzero blocks is a different requirement, because complete
layers can cancel. Without a separate nontermination or exact-termination
argument, enumeration alone does not prove that another nonzero block exists.

\subsection{Newton iteration with growing precision}

Write $E(U)=0$ for \eqref{eq:normalized}, and let
$\delta_* = \min_i\delta_i$. The zero jet is correct below $\delta_*$.
If $U_s$ is correct strictly below weight $s$, one Newton step performed
strictly below
\[
 h=\min(2s,H),\qquad
 U_h=\left[U_s-\frac{E(U_s)}{E'(U_s)}\right]_{<h},
\]
is correct strictly below $h$. Indeed, $E'(U_s)$ has nonzero constant
coefficient $p$, and the Newton error is quadratic in the old error. Products,
unit powers, and logarithmic composition do not decrease the weight of the
error. The argument uses exponent valuation, with polynomial logarithms as
coefficients; logarithmic degree need not stay bounded as $h$ increases.

The working cutoffs are
$\min(2\delta_*,H),\min(4\delta_*,H),\ldots,H$, rather than $H$ at every
step. Their construction uses exact comparisons, with no floating-point
estimate of the required iteration count. The identity $[E(U_h)]_{<h}=0$ characterizes the formal accuracy of each step. Together with triangular uniqueness it proves that the retained coefficients agree with those of the inverse. A numerical bound at a particular nonzero $z$ additionally requires an analytic remainder estimate.

If a unit jet is already correct below $s$, refinement to $H>s$ starts with
that same jet and the same doubling argument. Truncation to a smaller
cutoff preserves correctness. For a pure monomial forward function the
unit correction is exactly zero, so no refinement is needed.

\subsection{Combine resonances before applying Euler operators}

Put
\[
 H_1(z,L)=\sum_i z^{\delta_i}B_i(L),\qquad
 H_1(z,L)^n=\sum_\sigma z^\sigma A_{n,\sigma}(L).
\]
The multinomial theorem gives
\[
 A_{n,\sigma}=
 \sum_{\substack{|\bk|=n\\\bk\cdot\bdelta=\sigma}}
 \frac{n!}{\bk!}\,B^{\bk}.
\]
At fixed depth $n$ and weight $\sigma$, all summands of
\eqref{eq:master} have the same differential operator. By linearity, the unit
bracket of $g^r$ can therefore be computed as
\begin{equation}
 \frac{g(y)^r}{z^r}
 =1+\sum_{n\ge1}\sum_\sigma z^\sigma
 \frac{(-1)^n r}{p^n n!}
 \prod_{j=1}^{n-1}(\D+r+\sigma+pj)A_{n,\sigma}(L).
 \label{eq:grouped-lagrange}
\end{equation}
For cutoff $H$, successive sparse products compute
$[H_1^n]_{<H}$. Only $n$ with $n\delta_*<H$ can contribute, so the computation
is finite. Coefficients of equal weight are combined during each sparse
product, before differentiation. Contributions at the same weight but
different depths are then combined in the result; they too can cancel.

Equation~\eqref{eq:grouped-lagrange} is especially useful when the generators
have many additive relations. It does not require commensurable gaps, and it
does not replace the individual-coefficient formula: that formula
retains the distinct perturbation markers and their provenance. Grouping replaces several applications of the same linear differential
operator by one application to the sum. The boundary argument of
Section~\ref{sec:remainders} is still needed to identify a complete omitted
weight. Neither grouping nor marker-depth truncation changes the ordering
of the final nonzero blocks. The equivalence of the grouped formula,
the individual multi-index formula and Newton iteration follows from
linearity and formal uniqueness.
